\documentclass[
    language=english,
    doctype=article,
    institution=none,
]{../../omnilatex}

\RequirePackage{config/document-settings}
\addbibresource{../../bib/bibliography.bib}

\begin{document}

\maketitle

\begin{abstract}
DataPipe is an open-source streaming ETL (Extract, Transform, Load) platform
for real-time data integration.  This guide covers architecture, installation,
configuration, APIs, troubleshooting, and version history.
\end{abstract}

\tableofcontents

%=============================================================================
\section{Overview}
\label{sec:overview}
%=============================================================================

DataPipe ingests data from heterogeneous sources, applies configurable
transformation pipelines, and delivers processed records to downstream systems
with sub-second latency on Apache Kafka and Apache Flink.

Key characteristics:
\begin{itemize}
    \item \textbf{Throughput:} 500k events/s on a three-node cluster.
    \item \textbf{Latency:} End-to-end under 200\,ms at P95.
    \item \textbf{Durability:} At-least-once; exactly-once for supported sinks.
    \item \textbf{Extensibility:} Java/Scala plugin API for custom connectors.
\end{itemize}

\subsection{Architecture Diagram}

\begin{figure}[htbp]
    \centering
    \begin{tikzpicture}[
        node distance=1.2cm and 1.8cm,
        box/.style={draw, rounded corners, minimum width=2.8cm,
                     minimum height=0.9cm, text centered, font=\small,
                     fill=#1!20, draw=#1!60!black},
        arr/.style={->, >=stealth, thick},
    ]
        \node[box=blue]   (src)   {Data Sources};
        \node[box=teal, right=of src]    (ingest)  {Ingestion};
        \node[box=orange, right=of ingest] (kafka) {Kafka};
        \node[box=red, below=of kafka]   (flink)   {Flink};
        \node[box=purple, below=of flink] (sink)   {Sinks};
        \node[box=cyan, left=of flink]   (web)    {Web UI};
        \node[box=green!60!black, left=of sink] (pg) {PostgreSQL};
        \draw[arr] (src)--(ingest); \draw[arr] (ingest)--(kafka);
        \draw[arr] (kafka)--(flink); \draw[arr] (flink)--(sink);
        \draw[arr] (web)--(flink);   \draw[arr] (flink)--(pg);
    \end{tikzpicture}
    \caption{High-level architecture.}
    \label{fig:arch}
\end{figure}

%=============================================================================
\section{System Requirements}
\label{sec:requirements}
%=============================================================================

\begin{table}[htbp]
    \centering
    \caption{Hardware specifications.}
    \label{tab:hardware}
    \begin{tabular}{@{}lcc@{}}
        \toprule
        \textbf{Resource} & \textbf{Minimum} & \textbf{Recommended} \\
        \midrule
        CPU cores    & 4   & 16   \\
        RAM          & 8\,GB & 64\,GB \\
        Disk         & 50\,GB SSD & 500\,GB NVMe \\
        Network      & 1\,Gbps & 10\,Gbps \\
        \bottomrule
    \end{tabular}
\end{table}

Software: Java 17+, Kafka 3.6+, PostgreSQL 15+, Docker 24+ with Compose v2.

%=============================================================================
\section{Installation}
\label{sec:installation}
%=============================================================================

\begin{table}[htbp]
    \centering
    \caption{Platform-specific installation methods.}
    \label{tab:platforms}
    \begin{tabular}{@{}lp{5cm}p{3.5cm}@{}}
        \toprule
        \textbf{Platform} & \textbf{Method} & \textbf{Notes} \\
        \midrule
        Linux (x86\_64)  & Package manager / Docker & \texttt{.deb} and \texttt{.rpm} \\
        Linux (ARM64)    & Docker & Multi-arch image \\
        macOS (ARM)      & Homebrew & \texttt{brew install datapipeline/datapipe/datapipe} \\
        Windows          & Docker Desktop & WSL 2 required \\
        Kubernetes       & Helm chart & Production recommended \\
        \bottomrule
    \end{tabular}
\end{table}

\subsection{Quick Start}

\begin{minted}[caption={Docker Compose quick start.}, label={lst:docker-qs}]{bash}
git clone https://github.com/datapipeline/datapipe.git
cd datapipe && docker compose -f docker-compose.dev.yaml up -d
\end{minted}

\subsection{Kubernetes Production}

\begin{minted}[caption={Helm install.}, label={lst:helm}]{bash}
helm repo add datapipe https://charts.datapipeline.io/stable && helm repo update
helm install datapipe datapipe/datapipe -n datapipe --create-namespace -f values-prod.yaml
\end{minted}

%=============================================================================
\section{Configuration}
\label{sec:configuration}
%=============================================================================

\begin{minted}[
    caption={\texttt{datapipe.yaml} configuration file.},
    label={lst:config}, linenos,
]{yaml}
kafka:
  brokers: "kafka-1:9092,kafka-2:9092,kafka-3:9092"
  group_id: "datapipe-prod"
  session_timeout_ms: 30000
flink:
  parallelism: 16
  checkpoint_interval: 30000
  state_backend: "rocksdb"
pg:
  url: "jdbc:postgresql://db.internal:5432/datapipe"
  username: "${DB_USER}"
  password: "${DB_PASSWORD}"
  pool_size: 20
logging: { level: "INFO", format: "json" }
server: { host: "0.0.0.0", port: 8080 }
\end{minted}

\begin{table}[htbp]
    \centering
    \caption{Environment variables (override YAML values).}
    \label{tab:env-vars}
    \begin{tabular}{@{}lp{5.5cm}p{3cm}@{}}
        \toprule
        \textbf{Variable} & \textbf{Description} & \textbf{Default} \\
        \midrule
        \texttt{DATAPIPE\_KAFKA\_BROKERS} & Broker addresses & \texttt{localhost:9092} \\
        \texttt{DATAPIPE\_KAFKA\_GROUP\_ID} & Consumer group & \texttt{datapipe-default} \\
        \texttt{DATAPIPE\_FLINK\_PARALLELISM} & Parallelism & \texttt{4} \\
        \texttt{DATAPIPE\_PG\_URL} & JDBC URL & --- \\
        \texttt{DATAPIPE\_PG\_USER} & DB username & --- \\
        \texttt{DATAPIPE\_PG\_PASSWORD} & DB password & --- \\
        \texttt{DATAPIPE\_LOG\_LEVEL} & Verbosity & \texttt{INFO} \\
        \texttt{DATAPIPE\_SERVER\_PORT} & Listen port & \texttt{8080} \\
        \bottomrule
    \end{tabular}
\end{table}

\begin{table}[htbp]
    \centering
    \caption{Feature flags.}
    \label{tab:flags}
    \begin{tabular}{@{}lp{5cm}cc@{}}
        \toprule
        \textbf{Flag} & \textbf{Description} & \textbf{Default} & \textbf{Since} \\
        \midrule
        \texttt{exactly\_once\_sinks} & Exactly-once for sinks & \texttt{false} & v2.4.0 \\
        \texttt{schema\_registry} & Confluent Registry auto-register & \texttt{false} & v2.3.0 \\
        \texttt{grpc\_ingestion} & gRPC ingestion endpoint & \texttt{true} & v2.3.0 \\
        \texttt{rest\_ingestion} & Legacy REST ingestion & \texttt{false} & v2.3.0 \\
        \texttt{metrics\_export} & Prometheus \texttt{/metrics} & \texttt{true} & v2.1.0 \\
        \texttt{saml\_auth} & SAML 2.0 SSO & \texttt{false} & v2.2.0 \\
        \bottomrule
    \end{tabular}
\end{table}

%=============================================================================
\section{API Reference}
\label{sec:api-reference}
%=============================================================================

\begin{table}[htbp]
    \centering
    \caption{REST API endpoints.}
    \label{tab:rest-api}
    \begin{tabular}{@{}llp{4cm}p{2.5cm}@{}}
        \toprule
        \textbf{Method} & \textbf{Endpoint} & \textbf{Description} & \textbf{Parameters} \\
        \midrule
        \texttt{GET}    & \texttt{/api/v1/pipelines} & List pipelines & \texttt{?status=active} \\
        \texttt{POST}   & \texttt{/api/v1/pipelines} & Create pipeline & JSON body \\
        \texttt{GET}    & \texttt{/api/v1/pipelines/\{id\}} & Get details & Path \texttt{id} \\
        \texttt{PUT}    & \texttt{/api/v1/pipelines/\{id\}} & Update pipeline & Path \texttt{id} \\
        \texttt{DELETE} & \texttt{/api/v1/pipelines/\{id\}} & Delete pipeline & Path \texttt{id} \\
        \texttt{POST}   & \texttt{/api/v1/pipelines/\{id\}/start} & Start & Path \texttt{id} \\
        \texttt{POST}   & \texttt{/api/v1/pipelines/\{id\}/stop} & Stop & Path \texttt{id} \\
        \texttt{GET}    & \texttt{/api/v1/metrics} & Metrics & \texttt{?window=5m} \\
        \bottomrule
    \end{tabular}
\end{table}

\begin{minted}[
    caption={Create pipeline --- request and response.},
    label={lst:create-pipeline},
]{json}
// POST /api/v1/pipelines
{ "name": "user-events-etl",
  "source": { "type": "kafka", "topic": "raw-events" },
  "sink":   { "type": "postgresql", "table": "events" } }
// Response 201 Created
{ "id": "pl_8f3a2b1c", "status": "created", "created_at": "2026-06-12T10:30:00Z" }
\end{minted}

\subsection{Authentication}

\begin{minted}[caption={Obtain and use a Bearer token.}, label={lst:auth}]{bash}
curl -X POST https://auth.datapipeline.io/oauth/token \
    -d 'grant_type=client_credentials' -d 'client_id=ID' -d 'client_secret=SECRET'
# Then use the token:
curl -H "Authorization: Bearer <token>" https://api.datapipeline.io/api/v1/pipelines
\end{minted}

\begin{table}[htbp]
    \centering
    \caption{Error codes.}
    \label{tab:errors}
    \begin{tabular}{@{}clp{6.5cm}@{}}
        \toprule
        \textbf{Code} & \textbf{HTTP} & \textbf{Description} \\
        \midrule
        \texttt{PIPE\_NOT\_FOUND}      & 404 & Pipeline not found \\
        \texttt{PIPE\_ALREADY\_RUNNING} & 409 & Pipeline already running \\
        \texttt{INVALID\_CONFIG}       & 422 & Invalid configuration \\
        \texttt{AUTH\_REQUIRED}        & 401 & Missing Authorization header \\
        \texttt{AUTH\_EXPIRED}         & 401 & Token expired \\
        \texttt{FORBIDDEN}             & 403 & Insufficient permissions \\
        \texttt{RATE\_LIMITED}         & 429 & Too many requests \\
        \texttt{INTERNAL\_ERROR}       & 500 & Server error \\
        \bottomrule
    \end{tabular}
\end{table}

%=============================================================================
\section{CLI Reference}
\label{sec:cli-reference}
%=============================================================================

\begin{table}[htbp]
    \centering
    \caption{CLI commands.}
    \label{tab:cli}
    \begin{tabular}{@{}lp{5cm}p{3.5cm}@{}}
        \toprule
        \textbf{Command} & \textbf{Description} & \textbf{Example} \\
        \midrule
        \texttt{init}     & Scaffold a project & \texttt{datapipe init my-pipe} \\
        \texttt{validate} & Validate config & \texttt{datapipe validate pipeline.yaml} \\
        \texttt{run}      & Execute locally & \texttt{datapipe run --dry-run} \\
        \texttt{deploy}   & Deploy remote & \texttt{datapipe deploy --env prod} \\
        \texttt{status}   & Pipeline status & \texttt{datapipe status pl\_8f3a2b1c} \\
        \texttt{logs}     & Tail logs & \texttt{datapipe logs pl\_8f3a2b1c -f} \\
        \texttt{version}  & Print version & \texttt{datapipe version} \\
        \bottomrule
    \end{tabular}
\end{table}

%=============================================================================
\section{Code Examples}
\label{sec:code-examples}
%=============================================================================

\subsection{Python}

\begin{minted}[caption={Python client.}, label={lst:python}, language=Python, linenos]{python}
import requests

class DataPipeClient:
    def __init__(self, base_url: str, token: str):
        self.base = base_url.rstrip("/")
        self.s = requests.Session()
        self.s.headers["Authorization"] = f"Bearer {token}"

    def list_pipelines(self, status=None):
        r = self.s.get(f"{self.base}/api/v1/pipelines",
                       params={"status": status} if status else {})
        r.raise_for_status()
        return r.json()["pipelines"]

    def create_pipeline(self, config):
        r = self.s.post(f"{self.base}/api/v1/pipelines", json=config)
        r.raise_for_status()
        return r.json()

client = DataPipeClient("https://api.datapipeline.io", "YOUR_TOKEN")
for p in client.list_pipelines("active"):
    print(p["name"], p["status"])
\end{minted}

\subsection{Rust}

\begin{minted}[caption={Rust client.}, label={lst:rust}, language=Rust, linenos]{rust}
use serde::{Deserialize, Serialize};
use reqwest::Client;

#[derive(Debug, Serialize, Deserialize)]
struct Pipeline { id: String, name: String, status: String }

struct DataPipeClient { client: Client, base: String, token: String }

impl DataPipeClient {
    fn new(url: &str, tok: &str) -> Self {
        Self { client: Client::new(), base: url.trim_end_matches('/').into(), token: tok.into() }
    }
    async fn list_pipelines(&self) -> Result<Vec<Pipeline>, reqwest::Error> {
        let v = self.client.get(format!("{}/api/v1/pipelines", self.base))
            .bearer_auth(&self.token).send().await?
            .json::<serde_json::Value>().await?;
        Ok(serde_json::from_value(v["pipelines"].clone()).unwrap())
    }
}
\end{minted}

\subsection{Shell}

\begin{minted}[caption={Health check script.}, label={lst:shell}, language=bash, linenos]{bash}
#!/usr/bin/env bash
set -euo pipefail
API="${DATAPIPE_API_URL:-https://api.datapipeline.io}"
TOKEN="${DATAPIPE_TOKEN:?Missing}" && for i in 1 2 3; do
    [[ $(curl -s -o /dev/null -w "%{http_code}" -H "Authorization: Bearer ${TOKEN}" "${API}/health/ready") == 200 ]] && exit 0
    sleep 5; done; echo "Unhealthy" >&2; exit 1
\end{minted}

%=============================================================================
\section{Troubleshooting}
\label{sec:troubleshooting}
%=============================================================================

\begin{description}
    \item[Kafka connection refused] Check brokers and \texttt{kafka.brokers} config.
    \item[High checkpoint duration] Increase \texttt{flink.checkpoint\_timeout} or tune windows.
    \item[Out of memory] Set \texttt{JAVA\_OPTS=-Xmx16g}; monitor GC logs.
    \item[Schema evolution failures] Update pipeline definition before restarting.
    \item[401 auth errors] Token may be expired (default 3600\,s); re-authenticate.
    \item[Pipeline stuck starting] Check Flink logs for backpressure; increase parallelism.
\end{description}

Health endpoints: \texttt{/health/ready}, \texttt{/health/live}, \texttt{/metrics} (Prometheus).

%=============================================================================
\section{Architecture}
\label{sec:architecture}
%=============================================================================

\begin{figure}[htbp]
    \centering
    \begin{tikzpicture}[
        node distance=1.0cm and 1.5cm,
        comp/.style={draw, rounded corners, minimum width=2.4cm,
                     minimum height=0.8cm, text centered, font=\small,
                     fill=#1!15, draw=#1!50!black},
        arr/.style={->, >=stealth, thick},
    ]
        \node[comp=orange]              (api)   {REST API};
        \node[comp=teal,  right=of api]         (engine){Stream Engine};
        \node[comp=blue,  right=of engine]      (sink)  {Sink Manager};
        \node[comp=purple,below=of engine]      (auth)  {Auth Service};
        \node[comp=green!60!black,below=of api] (cfg)   {Config Store};
        \node[comp=red,   below=of sink]        (mon)   {Metrics};
        \draw[arr](api)--(engine); \draw[arr](engine)--(sink);
        \draw[arr](auth)--(api);   \draw[arr](cfg)--(api);
        \draw[arr](mon)--(engine);
    \end{tikzpicture}
    \caption{Component diagram.}
    \label{fig:components}
\end{figure}

\begin{figure}[htbp]
    \centering
    \begin{tikzpicture}[
        node distance=1.5cm and 2.0cm,
        stage/.style={draw, rounded corners, minimum width=2.2cm,
                      minimum height=0.8cm, text centered, font=\small},
        arr/.style={->, >=stealth, thick},
    ]
        \node[stage, fill=blue!20]     (src)   {Source};
        \node[stage, fill=teal!20, right=of src]   (parse) {Parse};
        \node[stage, fill=orange!20, right=of parse] (xform){Transform};
        \node[stage, fill=red!20, right=of xform]   (valid){Validate};
        \node[stage, fill=purple!20, right=of valid] (sink) {Sink};
        \node[stage, fill=gray!20, below=of xform]  (dlq)  {Dead-Letter Queue};
        \draw[arr](src)--(parse);
        \draw[arr](parse)--node[above,font=\scriptsize]{JSON/Avro}(xform);
        \draw[arr](xform)--node[above,font=\scriptsize]{Enriched}(valid);
        \draw[arr](valid)--node[above,font=\scriptsize]{Verified}(sink);
        \draw[arr,dashed](valid)--node[right,font=\scriptsize]{invalid}(dlq);
    \end{tikzpicture}
    \caption{Data flow diagram.}
    \label{fig:data-flow}
\end{figure}

\begin{figure}[htbp]
    \centering
    \begin{tikzpicture}[
        node distance=1.2cm and 1.5cm,
        zone/.style={draw,dashed,rounded corners,inner sep=10pt,label={[font=\small]above:#1}},
        box/.style={draw,rounded corners,min width=2.0cm,min height=0.7cm,
                     text centered,font=\small,fill=#1!15,draw=#1!50!black},
        arr/.style={->,>=stealth},
    ]
        \begin{scope}[zone=Kubernetes]
            \node[box=blue]  (ingress){Ingress};
            \node[box=teal, right=of ingress] (web) {Web UI};
            \node[box=orange,right=of web]    (api2){API};
            \node[box=red,  below=of api2]    (fl)  {Flink};
        \end{scope}
        \begin{scope}[zone=External]
            \node[box=purple,below=of fl]  (ka){Kafka};
            \node[box=green!60!black,right=of ka] (pg){PostgreSQL};
        \end{scope}
        \draw[arr](ingress)--(web); \draw[arr](web)--(api2);
        \draw[arr](api2)--(fl);     \draw[arr](fl)--(ka);
        \draw[arr](fl)--(pg);
    \end{tikzpicture}
    \caption{Deployment topology.}
    \label{fig:deployment}
\end{figure}

%=============================================================================
\section{Changelog}
\label{sec:changelog}
%=============================================================================

\subsection*{v2.4.0 --- 2026-05-15}
\begin{itemize}
    \item Exactly-once semantics for PostgreSQL and MySQL sinks.
    \item Built-in data quality validation framework.
    \item 40\% faster Flink job recovery via incremental checkpointing.
    \item Fixed memory leak in Kafka consumer poll loop at high throughput.
\end{itemize}

\subsection*{v2.3.0 --- 2026-03-01}
\begin{itemize}
    \item gRPC ingestion API for high-throughput producers.
    \item Apache Avro schema registry integration.
    \item Default serialization changed from JSON to Protocol Buffers.
    \item Deprecated legacy REST ingestion (\texttt{/api/v1/ingest}).
\end{itemize}

\subsection*{v2.2.0 --- 2026-01-10}
\begin{itemize}
    \item Helm chart for Kubernetes deployments.
    \item SAML 2.0 SSO for the web UI.
    \item Dropped Java 11 support (Java 17+ required).
    \item Fixed race condition in pipeline status reporting.
    \item Fixed timezone-aware timestamp handling.
\end{itemize}

\printbibliography[heading=none]

\end{document}
